\chapter{Estimation of the \texorpdfstring{$\ell_2$}{l2}-norm of the reward vector}
\label{chap:norm}

This chapter formalizes Section~3 and Appendix~D of the paper (the four subsections on
$\ell_2$-norm estimation): on the action set $\Xs = \ball_d$ one wants to estimate
$r := \norm{\theta}$ up to an additive error $\eps$ with probability at least $1-\delta$,
using $O(d\log(1/\delta)/\eps^2)$ observations. The target is Theorem~8 of the paper
(Theorem~\ref{thm:norm_estimation}). The procedure has three ingredients: a statistic
built from random Rademacher directions, each sampled repeatedly, which estimates $r^2$ with
an additive error proportional to $r\eps$ once a constant-factor estimate $r_0$ of $r$ is
available (Section~\ref{sec:norm_additive}); an adaptive multi-scale test which produces
such an $r_0$, or detects that $r$ is either at most $\eps$ or at least $\sqrt d$
(Section~\ref{sec:norm_test}); and a direct estimator for the large-norm regime
$r \ge \sqrt d$ (Section~\ref{sec:norm_large}). The three are composed into a
meta-algorithm in Section~\ref{sec:norm_meta}. The chapter is fully formalized.

From Chapter~\ref{chap:pre_subexp} we use the sub-exponential calculus: the definition
\ref{def:subexp}, the tail bound for averages (Corollary~\ref{cor:subexp_average_tail}), the
sub-exponential parameters of squared Gaussians (Lemma~\ref{lem:gaussian_square_subexp}), of
chi-square variables (Lemma~\ref{lem:chi_square_subexp}) and of squared Rademacher linear
forms (Lemma~\ref{lem:rademacher_square_subexp}). From the model chapter we use the seed
representation of the runs of a seeded algorithm (Lemma~\ref{lem:seeded}): the internal
randomness of the algorithm (the Rademacher directions) and the Gaussian noises are the
coordinates of an i.i.d.\ sequence, and every statistic of the run is a measurable
function of finitely many of these coordinates. Compared with the paper, the constants are
all explicit (no attempt was made to optimize them), the Hanson--Wright inequality is
replaced by Lemma~\ref{lem:rademacher_square_subexp}, the conditional argument for the noise
term is a Fubini argument over the directions, the large-norm regime uses a deterministic
basis design instead of random Rademacher rows, which removes the random-matrix
concentration step of the paper, and the sequential composition of the phases is replaced
by a deterministic round schedule on the seed space. Throughout the chapter $\eps \in (0,1]$, $\delta \in (0,1)$, $d \ge 1$,
$\Xs = \ball_d$ and $r := \norm{\theta}$.

\section{Estimation algorithms and the Rademacher statistic}
\label{sec:norm_setup}

\begin{definition}[Norm estimation algorithm]\label{def:norm_estimator}
\lean{Maiti2026Power.IsAccurateNormEst}\leanok
\uses{def:env, def:bai_alg}
An \emph{estimation algorithm with budget $T$} on $\ball_d$ is a pair $(\mathrm{alg}, \rho)$
of an LML algorithm with actions in $\ball_d$ and observations in $\R$, and a Markov kernel
$\rho$ from histories of length $T$ to $[0,\infty)$. Run against $\nu_\theta$
(Definition~\ref{def:env}) it produces $\hist_T$ and the estimate $\hat r \sim \rho(\hist_T)$;
we write $\Prob_\theta$ for the law of $(\hist_T, \hat r)$. It is
\emph{$(\eps,\delta)$-accurate} if for every $\theta \in \R^d$,
\[
  \Prob_\theta\big( \abs{\hat r - \norm{\theta}} \le \eps \big) \ge 1 - \delta .
\]
\end{definition}

\textbf{Lean remark.} This is Definition~\ref{def:bai_alg} with the recommendation type
$\Xs$ replaced by $[0,\infty)$ (or simply $\R$): the LML notion of an algorithm with a
recommendation is polymorphic in the recommendation type, so no new framework is needed.
All the estimators below are deterministic functions of the history, so $\rho$ is
\texttt{Kernel.deterministic}.

\begin{definition}[Rademacher direction]\label{def:rademacher_direction}
\lean{Maiti2026Power.radDir, Maiti2026Power.radDir_mem_unitBall, Maiti2026Power.measurable_radDir}\leanok
\uses{def:arm_set}
A \emph{Rademacher direction} is the random vector $x := \sigma/\sqrt d \in \R^d$ where
$\sigma = (\sigma_1, \dots, \sigma_d)$ has i.i.d.\ coordinates uniform on $\{-1, +1\}$. We
write $\mathrm{Rad}_d$ for its law (the image of the uniform measure on $\{-1,+1\}^d$ by
$\sigma \mapsto \sigma/\sqrt d$). Since $\norm{x} = 1$, $x \in \sphere^{d-1} \subset \ball_d$.
\end{definition}

\textbf{Lean remark.} The sign vector is encoded by a Boolean vector $u \in \{0,1\}^d$ with
the uniform law (\texttt{uniformBoolVec}, the seed of the meta-algorithm), through
$\sigma_i := \mathrm{signOf}(u_i) \in \{-1,+1\}$; $\mathrm{radDir}(u) := \sigma/\sqrt d$.

\begin{lemma}[Inner products with a Rademacher direction]\label{lem:direction_second_moment}
\lean{Maiti2026Power.inner_radDir, Maiti2026Power.card_mul_inner_radDir_sq, Maiti2026Power.norm_radDir_le_one}\leanok
\uses{def:rademacher_direction}
Let $x = \sigma/\sqrt d$ be a Rademacher direction and $\theta \in \R^d$. Then
$\ip{x}{\theta} = \ip{\sigma}{\theta}/\sqrt d$ and $d\,\ip{x}{\theta}^2 = \ip{\sigma}{\theta}^2$;
in particular $\E\big[\ip{x}{\theta}^2\big] = \norm{\theta}^2/d$.
\end{lemma}

\begin{proof}
\leanok
\uses{lem:rademacher_fourth_moment}
$\E[\sigma_i\sigma_j] = \indic\{i = j\}$ by independence and $\E\sigma_i = 0$,
$\sigma_i^2 = 1$, so $\E[\sigma\sigma^\top] = I_d$ and $\E[xx^\top] = I_d/d$. Then
$\E\ip{x}{\theta}^2 = \theta^\top\E[xx^\top]\theta = \norm{\theta}^2/d$ (this is also the
second-moment identity of Lemma~\ref{lem:rademacher_fourth_moment} divided by $d$). The last
identity is $\ip{x}{\theta} = \ip{\sigma}{\theta}/\sqrt d$.
\end{proof}

\begin{definition}[Rademacher block sampling and the squared statistic]\label{def:squared_statistic}
\lean{Maiti2026Power.rbMeasure, Maiti2026Power.rbMean, Maiti2026Power.rbStat, Maiti2026Power.rbX, Maiti2026Power.rbW, Maiti2026Power.rbStat_sub_sq_eq}\leanok
\uses{def:rademacher_direction, def:env, def:regret}
Let $s \ge 1$ and $K \ge 1$ be integers. The \emph{Rademacher block sampling rule}
$\mathrm{RB}(s, K)$ is the sampling rule with actions in $\ball_d$ and budget $sK$ which, for
each $k < K$, draws at round $ks$ a fresh Rademacher direction $x^{(k)} = \sigma^{(k)}/\sqrt d$
and plays the same action $x^{(k)}$ at rounds $ks + \ell$, $0 \le \ell < s$. Its
\emph{block space} is the product space
\[
  \Omega_{s,K} := \big(\{-1,+1\}^d\big)^K \times \big(\R^s\big)^K
  \quad\text{with the product law}\quad
  \Prob_{s,K} := \mathrm{Unif}\big(\{-1,+1\}^d\big)^{\otimes K} \otimes
    \big(\Normal(0,1)^{\otimes s}\big)^{\otimes K}
\]
(the directions of the $K$ blocks and the noises of the $sK$ rounds). For
$(\sigma, \eta) \in \Omega_{s,K}$ and a reward vector $\theta$ with $r = \norm{\theta}$, define
for $k < K$
\[
  x^{(k)} := \frac{\sigma^{(k)}}{\sqrt d}, \qquad
  \mu_k := \ip{x^{(k)}}{\theta}, \qquad
  \bar y_k := \mu_k + \frac1s \sum_{\ell < s} \eta_{k,\ell}, \qquad
  Z_k := d\Big( \bar y_k^2 - \frac1s \Big), \qquad
  \bar Z := \frac1K \sum_{k<K} Z_k ,
\]
\[
  X_k := d\mu_k^2 - r^2, \qquad
  W_k := Z_k - d\mu_k^2 .
\]
Then, as an algebraic identity,
\[
  \bar Z - r^2 = \frac1K \sum_{k<K} X_k + \frac1K \sum_{k<K} W_k .
\]
When $\mathrm{RB}(s,K)$ is run against $\nu_\theta$, the observations of block $k$ are
$y_{ks+\ell} = \mu_k + \eta_{k,\ell}$ (Lemma~\ref{lem:seeded}), so that $\bar y_k$ is the
average of the observations of block $k$ and $\bar Z$ is a statistic of the history.
\end{definition}

\textbf{Lean remark.} $\mathrm{RB}(s,K)$ is not a standalone algorithm in Lean: it appears as
the \emph{windows} of the meta-algorithm of Definition~\ref{def:meta_algorithm}, whose seeds
at the block starts are the sign vectors $\sigma^{(k)}$ and whose noises are the
$\eta_{k,\ell}$ (Lemma~\ref{lem:statistic_structure}). All the probabilistic statements of
this chapter about $\bar Z$ are statements about the block space
$(\Omega_{s,K}, \Prob_{s,K})$ (\texttt{rbMeasure}), and are transported to the run through the
window projections of Lemma~\ref{lem:statistic_structure}. The identity
$\bar Z - r^2 = \frac1K\sum X_k + \frac1K\sum W_k$ is a rearrangement of
$Z_k = X_k + W_k + r^2$.

\textbf{Lean remark.} The directions are the algorithm's internal randomness: in LML the
policy of $\mathrm{RB}(s,K)$ is the kernel which, at rounds $t$ with $s \mid t$, ignores the
history and returns the measure $\mathrm{Rad}_d$ (a finitely supported measure on the
subtype $\ball_d$), and at the other rounds returns the Dirac mass at $x_{t-1}$. The
quantities $\bar y_k$, $Z_k$, $\bar Z$ are measurable functions of the history of length
$sK$. The identity $\bar Z - r^2 = \frac1K\sum X_k + \frac1K\sum W_k$ is a rearrangement of
$Z_k = X_k + W_k + r^2$.

\begin{lemma}[Window projections of the seed space]\label{lem:statistic_structure}
\lean{Maiti2026Power.nsMeasure, Maiti2026Power.rbProj, Maiti2026Power.nsMeasure_map_rbProj, Maiti2026Power.indepFun_prefix_rbProj, Maiti2026Power.lnProj, Maiti2026Power.nsMeasure_map_lnProj, Maiti2026Power.indepFun_prefix_lnProj, ProbabilityTheory.indepFun_prefix_window, MeasureTheory.Measure.infinitePi_map_eval_comp, MeasureTheory.Measure.infinitePi_map_eval_comp', Maiti2026Power.shiftSeq, Maiti2026Power.nsMeasure_map_shiftSeq, Maiti2026Power.winStat, Maiti2026Power.winStat_eq_rbStat}\leanok
\uses{def:squared_statistic, lem:seeded}
Let $\Omega_{\mathrm{seed}} := \big(\{-1,+1\}^d \times \R\big)^{\N}$ with the product law
$\Prob_{\mathrm{seed}} := \big(\mathrm{Unif}(\{-1,+1\}^d) \otimes \Normal(0,1)\big)^{\otimes\N}$
(the seed space of Lemma~\ref{lem:seeded} for seeds which are uniform sign vectors and
standard Gaussian noises), and write $\omega_t = (\sigma_t, \eta_t)$ for its coordinates.
\begin{enumerate}
  \item[(i)] For $a \in \N$, $s, K \ge 1$, the \emph{Rademacher window projection}
    $\pi_{a,s,K}(\omega) := \big( (\sigma_{a+ks})_{k<K}, (\eta_{a+ks+\ell})_{k<K,\ell<s} \big)
    \in \Omega_{s,K}$ has law $\Prob_{s,K}$ (Definition~\ref{def:squared_statistic}), and it
    is independent of the prefix $(\omega_t)_{t<T}$ whenever $T \le a$.
  \item[(ii)] For $a, n \in \N$, the \emph{basis window projection}
    $\pi'_{a,n}(\omega) := (\eta_{a+in+\ell})_{i<d,\ell<n} \in (\R^n)^d$ has law
    $\Normal(0,1)^{\otimes dn}$, and it is independent of the prefix $(\omega_t)_{t<T}$
    whenever $T \le a$.
  \item[(iii)] For an observation sequence $y : \N \to \R$, let
    $\bar Z_{a,s,K}(y) := \frac1K\sum_{k<K} d\big( (\frac1s\sum_{\ell<s} y_{a+ks+\ell})^2 - \frac1s \big)$
    be the squared statistic of the window of rounds $[a, a + sK)$. If
    $y_{a+ks+\ell} = \ip{\sigma_{a+ks}/\sqrt d}{\theta} + \eta_{a+ks+\ell}$ for all $k < K$,
    $\ell < s$, then $\bar Z_{a,s,K}(y) = \bar Z\big(\pi_{a,s,K}(\omega)\big)$, the statistic of
    Definition~\ref{def:squared_statistic} evaluated at the window projection.
  \item[(iv)] For $a \in \N$, the shifted sequence $(\omega_{a+t})_{t\in\N}$ has law
    $\Prob_{\mathrm{seed}}$.
\end{enumerate}
\end{lemma}

\begin{proof}
\leanok
\uses{lem:seeded, def:squared_statistic}
The coordinates of an i.i.d.\ sequence indexed by an injective family of indices are
i.i.d.\ with the same law: the image of $\mu^{\otimes\N}$ by
$\omega \mapsto (\omega_{f(p)})_{p \in J}$ is $\mu^{\otimes J}$ for injective $f : J \to \N$
(\texttt{Measure.infinitePi\_map\_eval\_comp}, from the restriction and reindexing lemmas of
infinite product measures; for an infinite index set, as in (iv), the image measure is
identified through its finite-dimensional cylinders, \texttt{Measure.eq\_infinitePi}). The window projections are of this form (the indices
$a + ks + \ell$, resp.\ $a + in + \ell$, are pairwise distinct), which gives the laws after
splitting the pairs $(\sigma, \eta)$ and currying the index $(k, \ell)$. The prefix and a
window with $T \le a$ use disjoint sets of coordinates, and disjoint families of
coordinates of an i.i.d.\ sequence are independent (\texttt{iIndepFun\_infinitePi} and
\texttt{indepFun\_finset}). Part (iii) is a rewriting: the block average of the
observations is $\mu_k + \frac1s\sum_\ell \eta_{k,\ell}$.
\end{proof}

\textbf{Lean remark.} The lemma replaces the sequential composition and the conditional
independence arguments of the paper: on the seed space, every event of the multi-scale
phase (rounds $t < T_1$) is a function of the prefix $(\omega_t)_{t<T_1}$, and every
statistic of the second phase is a function of a window projection with $a = T_1$; the two
are independent by (i)--(ii).

\begin{proof}
\uses{lem:noise_representation, lem:gauss_pi, lem:gauss_linear_image}
By Lemma~\ref{lem:noise_representation} the $\eta_t$, $t < sK$, are i.i.d.\ $\Normal(0,1)$
and $\eta_t$ is independent of $(\hist_t, x_t)$. Since $x_{ks+\ell} = x^{(k)}$ for
$\ell < s$, averaging the observations of block $k$ gives $\bar y_k = \mu_k + \zeta_k$.
The direction $x^{(k)}$ is drawn at round $ks$ from the constant kernel $\mathrm{Rad}_d$, so
it is independent of $\hist_{ks}$, hence of $(x^{(j)}, \zeta_j)_{j<k}$; and
$(\eta_{ks+\ell})_{\ell<s}$ is independent of $(\hist_{ks}, x^{(k)})$. An induction on $k$
shows that the law of $\big((x^{(j)},\zeta_j)\big)_{j \le k}$ is the product law: at each
step the conditional law of $(x^{(k)}, (\eta_{ks+\ell})_{\ell<s})$ given the past is the
constant kernel $\mathrm{Rad}_d \otimes \Normal(0,1)^{\otimes s}$ (Lemma~\ref{lem:gauss_pi}).
Finally $\zeta_k = \frac1s\mathbf 1^\top(\eta_{ks+\ell})_{\ell<s} \sim \Normal(0, 1/s)$ by
Lemma~\ref{lem:gauss_linear_image}.
\end{proof}

\textbf{Lean remark.} The proof is an induction on the blocks using the step decomposition
of the trajectory measure (LML \texttt{IsAlgEnvSeq.hasCondDistrib\_step}): since the policy
kernel at round $ks$ does not depend on the history, the conditional law of the block
$k$ given $\hist_{ks}$ is a constant kernel, which is exactly independence from the past.
Alternatively one can view $\mathrm{RB}(s,K)$ as the $K$-fold sequential composition
(Lemma~\ref{lem:composition}) of the one-block rule $\mathrm{RB}(s,1)$.

\section{Additive estimation given a constant-factor estimate}
\label{sec:norm_additive}

Throughout this section $\mathrm{RB}(s,K)$ is run against $\nu_\theta$ and the notation of
Definition~\ref{def:squared_statistic} is in force. The decomposition
$\bar Z - r^2 = \frac1K\sum_k X_k + \frac1K\sum_k W_k$ separates the fluctuations of the
directions (Term~1, the $X_k$) from the fluctuations of the noise (Term~2, the $W_k$).

\begin{lemma}[Term 1 is sub-exponential]\label{lem:term1_subexp}
\lean{Maiti2026Power.hasSubexponentialMGF_rbX, Maiti2026Power.iIndepFun_rbX, Maiti2026Power.rbX_eq_zero_of_eq_zero}\leanok
\uses{def:squared_statistic, lem:statistic_structure, def:subexp}
On the block space $(\Omega_{s,K}, \Prob_{s,K})$, the variables $X_k = d\mu_k^2 - r^2$,
$k < K$, are independent and each has $\mathrm{HasSubexponentialMGF}\ X_k\ (16 r^4)\ (4r^2)$.
If $r = 0$ then $X_k = 0$ for all $k$.
\end{lemma}

\begin{proof}
\leanok
\uses{lem:rademacher_square_subexp, lem:direction_second_moment, def:squared_statistic}
By Lemma~\ref{lem:direction_second_moment}, $d\mu_k^2 = \ip{\sigma^{(k)}}{\theta}^2$ where
$\sigma^{(k)}$ is a uniform sign vector, so $X_k = \ip{\sigma^{(k)}}{\theta}^2 - \norm{\theta}^2$,
which has the sub-exponential parameters $(16r^4, 4r^2)$ by
Lemma~\ref{lem:rademacher_square_subexp}. The $\sigma^{(k)}$ are i.i.d.\ by
Lemma~\ref{lem:statistic_structure}, hence so are the $X_k$. When $\theta = 0$,
$\mu_k = 0 = r$.
\end{proof}

\begin{lemma}[Tail bound for Term 1]\label{lem:term1_bound}
\lean{Maiti2026Power.rbMeasure_real_avg_rbX_gt_le}\leanok
\uses{lem:term1_subexp}
Assume $r > 0$. For every $t \ge 0$,
\[
  \Prob_{s,K}\Big( \Big| \frac1K\sum_{k<K} X_k \Big| > t \Big)
  \le 2\exp\Big( -K \min\Big( \frac{t^2}{32 r^4}, \frac{t}{8 r^2} \Big) \Big).
\]
If $r = 0$ the left-hand side is $0$.
\end{lemma}

\begin{proof}
\leanok
\uses{cor:subexp_average_tail, lem:term1_subexp}
Corollary~\ref{cor:subexp_average_tail} applied to the independent variables $X_k$ with the
common parameters $(V, b) = (16r^4, 4r^2)$ (Lemma~\ref{lem:term1_subexp}) gives the bound
$2\exp(-K\min(t^2/(2V), t/(2b)))$, and $2V = 32r^4$, $2b = 8r^2$. The case $r = 0$ is
$X_k = 0$.
\end{proof}

\begin{lemma}[Term 2 is conditionally sub-exponential]\label{lem:term2_conditional}
\lean{Maiti2026Power.rbNoise, Maiti2026Power.hasLaw_sqrt_mul_rbMean, Maiti2026Power.hasSubexponentialMGF_rbW, Maiti2026Power.iIndepFun_rbW}\leanok
\uses{def:squared_statistic, lem:statistic_structure, def:subexp}
Fix deterministic sign vectors $\sigma^{(0)}, \dots, \sigma^{(K-1)}$, let
$x^{(k)} := \sigma^{(k)}/\sqrt d$, $\mu_k := \ip{x^{(k)}}{\theta}$, and let the noises
$(\eta_{k,\ell})_{k<K,\ell<s}$ be i.i.d.\ $\Normal(0,1)$ (the noise marginal
$\Prob^{\eta}_{s,K} := (\Normal(0,1)^{\otimes s})^{\otimes K}$ of the block space). Define
$\bar y_k$, $Z_k$ and $W_k := Z_k - d\mu_k^2$ as in Definition~\ref{def:squared_statistic}.
Then $\sqrt d\,\bar y_k \sim \Normal(\sqrt d\,\mu_k, d/s)$, the $W_k$ are independent and
$\mathrm{HasSubexponentialMGF}\ W_k\ V_k\ b$ with
\[
  V_k := 8\Big( \frac{d^2}{s^2} + \mu_k^2\,\frac{d^2}{s} \Big), \qquad b := \frac{4d}{s} .
\]
\end{lemma}

\begin{proof}
\leanok
\uses{lem:gaussian_square_subexp, lem:gauss_linear_image}
$\sqrt d\,\bar y_k \sim \Normal(\sqrt d\,\mu_k, d/s)$ (Lemma~\ref{lem:gauss_linear_image}),
and
\[
  W_k = d\bar y_k^2 - \frac ds - d\mu_k^2
      = \big(\sqrt d\,\bar y_k\big)^2 - \Big( (\sqrt d\,\mu_k)^2 + \frac ds \Big),
\]
which is $X^2 - (\mu^2 + \sigma^2)$ for $X \sim \Normal(\mu,\sigma^2)$ with
$\mu = \sqrt d\,\mu_k$, $\sigma^2 = d/s$. Lemma~\ref{lem:gaussian_square_subexp} gives the
parameters $(8(\sigma^4 + \mu^2\sigma^2), 4\sigma^2) = (8(d^2/s^2 + \mu_k^2 d^2/s), 4d/s)$.
Independence of the $W_k$ follows from that of the $\zeta_k$, each $W_k$ being a function
of $\zeta_k$ alone.
\end{proof}

\begin{lemma}[Tail bound for Term 2]\label{lem:term2_bound}
\lean{Maiti2026Power.rbVbar, Maiti2026Power.sum_rbVbar_eq, Maiti2026Power.rbNoise_real_avg_rbW_ge_le, Maiti2026Power.rbMeasure_real_avg_rbW_gt_le}\leanok
\uses{lem:term2_conditional, lem:statistic_structure}
On the block space, let
\[
  \bar V := \frac1K\sum_{k<K} V_k
  = 8\Big( \frac{d^2}{s^2} + \frac{d^2}{s}\cdot\frac1K\sum_{k<K}\mu_k^2 \Big),
  \qquad b := \frac{4d}{s} ,
\]
a function of the directions. Then for every $\tau \in \R$ and $t \ge 0$,
\[
  \Prob_{s,K}\Big( \Big| \frac1K\sum_{k<K} W_k \Big| > t \Big)
  \le \Prob_{s,K}(\bar V > \tau)
    + 2\exp\Big( -K\min\Big( \frac{t^2}{2\tau}, \frac{t}{2b} \Big) \Big).
\]
\end{lemma}

\begin{proof}
\leanok
\uses{lem:term2_conditional, lem:statistic_structure, cor:subexp_average_tail}
By Lemma~\ref{lem:statistic_structure} the joint law of the directions and the noise
averages is the product $\mathrm{Rad}_d^{\otimes K} \otimes \Normal(0,1/s)^{\otimes K}$, and
$\frac1K\sum_k W_k$ is a measurable function $F(x^{(0:K)}, \zeta_{0:K})$. Split the event on
$\{\bar V \le \tau\}$, which depends on the directions only, and apply Fubini:
\[
  \Prob_\theta\Big(\Big|\frac1K\sum_k W_k\Big| > t,\ \bar V \le \tau\Big)
  = \int \indic\{\bar V(x) \le \tau\}\,
    \Prob_\zeta\Big( \Big| \frac1K\sum_k W_k(x, \zeta) \Big| > t \Big)\,
    \mathrm{Rad}_d^{\otimes K}(dx) .
\]
For fixed directions $x$, Lemma~\ref{lem:term2_conditional} and
Corollary~\ref{cor:subexp_average_tail} bound the inner probability by
$2\exp(-K\min(t^2/(2\bar V(x)), t/(2b)))$, which is at most
$2\exp(-K\min(t^2/(2\tau), t/(2b)))$ when $\bar V(x) \le \tau$. Adding
$\Prob_\theta(\bar V > \tau)$ for the complementary event gives the claim.
\end{proof}

\textbf{Lean remark.} The block space is the product
$\mathrm{Unif}(\{-1,+1\}^d)^{\otimes K} \otimes (\Normal(0,1)^{\otimes s})^{\otimes K}$ and the
bound is proved by Fubini (\texttt{MeasureTheory.Measure.prod}) over the directions, with
the inner bound of Lemma~\ref{lem:term2_conditional} for each fixed direction vector. No
conditional expectation is needed.

\begin{lemma}[Bound on the conditional variance proxy]\label{lem:vbar_bound}
\lean{Maiti2026Power.rbMeasure_real_rbVbar_gt_le}\leanok
\uses{lem:term2_bound, lem:term1_bound}
With $\bar V$ as in Lemma~\ref{lem:term2_bound} and
\[
  \tau := 8\Big( \frac{d^2}{s^2} + \frac{3 d r^2}{2 s} \Big),
  \qquad\text{one has}\qquad
  \Prob_{s,K}(\bar V > \tau) \le 2\exp(-K/128) .
\]
\end{lemma}

\begin{proof}
\leanok
\uses{lem:term1_bound, def:squared_statistic}
Since $\frac1K\sum_k d\mu_k^2 = r^2 + \frac1K\sum_k X_k$, on the event
$\{|\frac1K\sum_k X_k| \le r^2/2\}$ we have $\frac1K\sum_k\mu_k^2 \le 3r^2/(2d)$ and hence
$\bar V \le 8(d^2/s^2 + (d^2/s)\cdot 3r^2/(2d)) = \tau$. If $r = 0$ this event is sure and
$\bar V = 8d^2/s^2 \le \tau$. If $r > 0$, Lemma~\ref{lem:term1_bound} with $t = r^2/2$ gives
$\min(t^2/(32r^4), t/(8r^2)) = \min(1/128, 1/16) = 1/128$, so the complementary event has
probability at most $2\exp(-K/128)$.
\end{proof}

\begin{definition}[Additive estimation algorithm; Algorithm~1]\label{def:additive_algorithm}
\lean{Maiti2026Power.addS, Maiti2026Power.addK}\leanok
\uses{def:squared_statistic, def:norm_estimator}
Let $\eps \in (0,1]$, $\delta \in (0,1)$ and $r_0 > 0$ (a scale parameter), and set
\[
  c_0 := 4, \qquad c_1 := 4096, \qquad
  s(r_0) := \Big\lceil \frac{c_0 d}{r_0^2} \Big\rceil, \qquad
  K(r_0) := \Big\lceil \frac{c_1 r_0^2 \log(4/\delta)}{\eps^2} \Big\rceil .
\]
\emph{Algorithm~1} with parameters $(\eps,\delta,r_0)$, written
$\mathrm{AddEst}(\eps,\delta,r_0)$, is the estimation algorithm
(Definition~\ref{def:norm_estimator}) with budget $s(r_0)K(r_0)$ whose sampling rule is the
Rademacher block sampling rule $\mathrm{RB}(s(r_0),K(r_0))$ of
Definition~\ref{def:squared_statistic} and whose estimate is the deterministic function of
the history
\[
  \hat r := \sqrt{\max(\bar Z, 0)} ,
\]
where $\bar Z$ is the squared statistic of Definition~\ref{def:squared_statistic}. When
$\eps$ and $\delta$ are fixed we write $s := s(r_0)$ and $K := K(r_0)$.
\end{definition}

\textbf{Lean remark.} Only the parameters $s(r_0), K(r_0)$ (natural numbers computed by
\texttt{Nat.ceil}) are defined on their own; the algorithm itself is the second phase of the
meta-algorithm of Definition~\ref{def:meta_algorithm} for the data-dependent scale
$r_0 = t_{j_*}$, and its estimate is a case of \texttt{NormEstParam.estimate}. The guarantee
below is stated on the block space.

\textbf{Lean remark.} The algorithm is a pair (LML algorithm, deterministic recommendation
kernel); its parameters $s, K$ are natural numbers computed from $(\eps,\delta,r_0)$ by
\texttt{Nat.ceil}. The meta-algorithm of Definition~\ref{def:meta_algorithm} instantiates it
with a data-dependent $r_0$ taking finitely many values.

\begin{theorem}[Additive estimation with a constant-factor estimate; Algorithm~1]\label{thm:additive_estimation}
\lean{Maiti2026Power.rbMeasure_real_addEst_gt_le, Maiti2026Power.rbMeasure_real_sqrt_addEst_gt_le, Maiti2026Power.NormEstParam.s'_mul_K'_le}\leanok
\uses{def:additive_algorithm, def:norm_estimator, def:multiscale_algorithm}
Let $\eps \in (0,1]$, $\delta \in (0,1)$, $r_0 \ge \eps$, and let $s = s(r_0)$, $K = K(r_0)$
be the parameters of Algorithm~1 (Definition~\ref{def:additive_algorithm}, with
$c_0 = 4$, $c_1 = 4096$). Then for every
$\theta$ with $r/2 < r_0 \le 2r$, on the block space $(\Omega_{s,K}, \Prob_{s,K})$,
\[
  \Prob_{s,K}\Big( \abs{\bar Z - r^2} \le \frac{r\eps}{2} \Big) \ge 1 - \frac{3\delta}{8},
  \qquad\text{hence}\qquad
  \Prob_{s,K}\Big( \abs{\hat r - r} \le \frac{\eps}{2} \Big) \ge 1 - \frac{3\delta}{8}
  \ge 1 - \frac{\delta}{2} .
\]
Moreover, for the scales $r_0 = t_j = 2^j\eps$ with $t_j < 2\sqrt d$ (the scales $j < J$ of
Definition~\ref{def:multiscale_algorithm}), the budget satisfies
\[
  sK \le C\,\frac{d\log(4/\delta)}{\eps^2}, \qquad C := 32776 .
\]
\end{theorem}

\begin{proof}
\leanok
\uses{lem:term1_bound, lem:term2_bound, lem:vbar_bound, def:squared_statistic, def:additive_algorithm, def:multiscale_algorithm}
Write $L := \log(4/\delta) \ge \log 4 > 1$ and note $\log(16/\delta) = \log 4 + L \le 2L$.
Since $\eps \le r_0 \le 2r$ we have $r > 0$; the hypotheses also give $r_0/2 < r < 2r_0$. By
Definition~\ref{def:squared_statistic}, $\{|\bar Z - r^2| > r\eps/2\}$ is contained in the
union of $E_1 := \{|\frac1K\sum X_k| > r\eps/4\}$ and $E_2 := \{|\frac1K\sum W_k| > r\eps/4\}$.

\emph{Term 1.} Lemma~\ref{lem:term1_bound} with $t = r\eps/4$: $t^2/(32r^4) = \eps^2/(512r^2)$
and $t/(8r^2) = \eps/(32r)$. Since $\eps \le 2r \le 16r$, the minimum is $\eps^2/(512r^2)$,
so $\Prob(E_1) \le 2\exp(-K\eps^2/(512r^2))$. This is at most $\delta/8$ as soon as
$K \ge 512r^2\log(16/\delta)/\eps^2$, which holds because $r^2 < 4r_0^2$ and
$\log(16/\delta) \le 2L$ give $512r^2\log(16/\delta)/\eps^2 \le 4096\,r_0^2L/\eps^2 \le K$.

\emph{Variance proxy.} Lemma~\ref{lem:vbar_bound}: $\Prob(\bar V > \tau) \le 2e^{-K/128} \le \delta/8$
because $K \ge 4096 r_0^2L/\eps^2 \ge 4096 L \ge 128\log(16/\delta)$ (using $r_0 \ge \eps$).

\emph{Term 2.} Since $s \ge 4d/r_0^2$ we have $d/s \le r_0^2/4$, hence $b = 4d/s \le r_0^2$ and,
with $r^2 < 4r_0^2$,
\[
  \tau = 8\Big( \frac{d^2}{s^2} + \frac{3dr^2}{2s} \Big)
  \le 8\Big( \frac{r_0^4}{16} + \frac{3 r^2 r_0^2}{8} \Big)
  \le 8\Big( \frac{r_0^4}{16} + \frac{3 r_0^4}{2} \Big) = \frac{25}{2}\,r_0^4 .
\]
With $t = r\eps/4$ and $r > r_0/2$:
$t^2/(2\tau) \ge (r^2\eps^2/16)/(25r_0^4) = r^2\eps^2/(400r_0^4) > \eps^2/(1600 r_0^2)$ and
$t/(2b) \ge r\eps/(8r_0^2) > \eps/(16r_0)$; since $\eps \le r_0$,
$\eps^2/(1600r_0^2) \le \eps/(16r_0)$, so the minimum in Lemma~\ref{lem:term2_bound} is at
least $\eps^2/(1600r_0^2)$ and
$\Prob(E_2) \le \Prob(\bar V > \tau) + 2\exp(-K\eps^2/(1600r_0^2))$. The last term is at
most $\delta/8$ as soon as $K \ge 1600r_0^2\log(16/\delta)/\eps^2$, which holds since
$1600r_0^2\log(16/\delta)/\eps^2 \le 3200\,r_0^2L/\eps^2 \le K$.

\emph{Union bound.} $\Prob(|\bar Z - r^2| > r\eps/2) \le \delta/8 + \delta/8 + \delta/8 = 3\delta/8$.

\emph{From $r^2$ to $r$.} On the event $|\bar Z - r^2| \le r\eps/2$, since $r > r_0/2 \ge \eps/2$,
$\bar Z \ge r(r - \eps/2) > 0$, so $\hat r = \sqrt{\bar Z}$ and
\[
  \abs{\hat r - r} = \frac{\abs{\bar Z - r^2}}{\sqrt{\bar Z} + r} \le \frac{r\eps/2}{r} = \frac{\eps}{2}.
\]

\emph{Budget.} For $r_0 = t_j < 2\sqrt d$, $4d/r_0^2 > 1$, so
$s \le 4d/r_0^2 + 1 \le 8d/r_0^2 = (8d/\eps^2)\,4^{-j}$; and
$K \le 4096\,r_0^2L/\eps^2 + 1 = 4096\cdot 4^j L + 1$. Hence
$sK \le (8d/\eps^2)(4096\,L + 4^{-j}) \le (8d/\eps^2)(4096\,L + 1) \le 8\cdot 4097\, dL/\eps^2
= 32776\, dL/\eps^2$, using $L \ge 1$.
\end{proof}

\begin{remark}[On the constants of Algorithm 1]\label{rem:norm_additive_constants}
The paper takes $c_0, c_1$ ``large enough''; the values $c_0 = 4$, $c_1 = 4096$ are what the
computation above produces with the sub-exponential parameters $(16r^4, 4r^2)$ of
Lemma~\ref{lem:rademacher_square_subexp} and the threshold $t = r\eps/4$ of the paper. The
proof shows that $|\hat r - r| \le \eps/2$, twice as much as needed; using the threshold
$t = r\eps/2$ instead would divide $c_1$ (and $C$) by four. The three failure events are
bounded by $\delta/8$ each, so the total failure probability is $3\delta/8$ rather than the
$\delta/2$ claimed by the paper; the extra room is not used.
\end{remark}

\section{A constant-factor estimate by a multi-scale test}
\label{sec:norm_test}

\begin{definition}[Scale test]\label{def:scale_test}
\lean{Maiti2026Power.testS, Maiti2026Power.testK}\leanok
\uses{def:squared_statistic}
For a scale $t > 0$ and a confidence level $\delta' \in (0,1)$, the test
$\mathrm{Test}(t,\delta')$ is the algorithm $\mathrm{RB}(s,K)$ with
\[
  s := \Big\lceil \frac{c_0 d}{t^2} \Big\rceil, \qquad
  K := \big\lceil c_1' \log(1/\delta') \big\rceil, \qquad
  c_0 := 4, \quad c_1' := 5120 ,
\]
followed by the binary output $\hat H := H_1$ if $U \ge \tfrac32 t^2$ and $\hat H := H_0$
otherwise, where $U := \bar Z$ is the statistic of Definition~\ref{def:squared_statistic}.
The hypothesis $H_0$ is ``$r \le t$'' and $H_1$ is ``$r \ge 2t$''; when $t < r < 2t$ neither
output counts as an error. Its budget is $sK$.
\end{definition}

\textbf{Lean remark.} As for Algorithm~1, only the parameters are defined on their own; the
tests are the windows of the first phase of the meta-algorithm
(Definition~\ref{def:multiscale_algorithm}), and the error bounds below are stated on the
block space.

\begin{lemma}[Error of the test under $H_0$]\label{lem:test_error_h0}
\lean{Maiti2026Power.testS_facts, Maiti2026Power.testK_facts, Maiti2026Power.six_mul_exp_neg_le, Maiti2026Power.rbMeasure_real_test_h0_le}\leanok
\uses{def:scale_test}
Let $\delta' \in (0,1/4]$ and $\theta$ with $r \le t$. Then, on the block space of
$\mathrm{RB}(s,K)$, $\Prob_{s,K}(U \ge \tfrac32 t^2) \le 3\delta'/4$: the test returns $H_1$
with probability at most $3\delta'/4$.
\end{lemma}

\begin{proof}
\leanok
\uses{lem:term1_bound, lem:term2_bound, lem:vbar_bound, def:squared_statistic}
Write $T_1 := \frac1K\sum_k X_k$, $T_2 := \frac1K\sum_k W_k$, so $U = r^2 + T_1 + T_2$. If
$|T_1| \le t^2/8$ and $|T_2| \le t^2/8$ then $U \le r^2 + t^2/4 \le 5t^2/4 < 3t^2/2$ and the
output is $H_0$. Hence the error event is contained in
$\{|T_1| > t^2/8\} \cup \{|T_2| > t^2/8\}$.

\emph{Term 1.} If $r = 0$ then $T_1 = 0$. Otherwise Lemma~\ref{lem:term1_bound} with
$a_0 = t^2/8$ gives $a_0^2/(32r^4) = t^4/(2048 r^4) \ge 1/2048$ and
$a_0/(8r^2) = t^2/(64r^2) \ge 1/64$ since $r \le t$; so
$\Prob(|T_1| > t^2/8) \le 2e^{-K/2048}$.

\emph{Term 2.} Since $s \ge 4d/t^2$: $d/s \le t^2/4$, $b = 4d/s \le t^2$ and
$\tau = 8(d^2/s^2 + 3dr^2/(2s)) \le 8(t^4/16 + 3t^4/8) = 7t^4/2$ (using $r \le t$). With
$\nu_0 = t^2/8$: $\nu_0^2/(2\tau) \ge (t^4/64)/(7t^4) = 1/448$ and
$\nu_0/(2b) \ge (t^2/8)/(2t^2) = 1/16$. Lemmas~\ref{lem:term2_bound} and~\ref{lem:vbar_bound}
give $\Prob(|T_2| > t^2/8) \le 2e^{-K/128} + 2e^{-K/448}$.

\emph{Total.} The error probability is at most
$2e^{-K/2048} + 2e^{-K/128} + 2e^{-K/448} \le 6e^{-K/2048}$. Since $\delta' \le 1/4$,
$\log 8 = \tfrac32\log 4 \le \tfrac32\log(1/\delta')$, hence
$2048\log(8/\delta') \le 2048\cdot\tfrac52\log(1/\delta') = 5120\log(1/\delta') \le K$, and
$6e^{-K/2048} \le 6\delta'/8 = 3\delta'/4$.
\end{proof}

\begin{lemma}[Error of the test under $H_1$]\label{lem:test_error_h1}
\lean{Maiti2026Power.rbMeasure_real_test_h1_le}\leanok
\uses{def:scale_test}
Let $\delta' \in (0,1/4]$ and $\theta$ with $r \ge 2t$. Then
$\Prob_{s,K}(U < \tfrac32 t^2) \le 3\delta'/4$: the test returns $H_0$ with probability at
most $3\delta'/4$.
\end{lemma}

\begin{proof}
\leanok
\uses{lem:term1_bound, lem:term2_bound, lem:vbar_bound, def:squared_statistic}
With $T_1, T_2$ as above, if $|T_1| \le r^2/8$ and $|T_2| \le r^2/8$ then
$U \ge r^2 - r^2/4 = 3r^2/4 \ge 3t^2$ (as $r^2 \ge 4t^2$) and the output is $H_1$. Hence the
error event is contained in $\{|T_1| > r^2/8\} \cup \{|T_2| > r^2/8\}$; here $r > 0$.

\emph{Term 1.} Lemma~\ref{lem:term1_bound} with $a_1 = r^2/8$:
$\min(a_1^2/(32r^4), a_1/(8r^2)) = \min(1/2048, 1/64) = 1/2048$, so
$\Prob(|T_1| > r^2/8) \le 2e^{-K/2048}$.

\emph{Term 2.} As before $d/s \le t^2/4 \le r^2/16$ and $b \le t^2 \le r^2/4$, and
$\tau \le 8(t^4/16 + 3t^2r^2/8) \le 8(r^4/256 + 3r^4/32) = 25r^4/32$. With $\nu_1 = r^2/8$:
$\nu_1^2/(2\tau) \ge (r^4/64)/(25r^4/16) = 1/100$ and $\nu_1/(2b) \ge (r^2/8)/(r^2/2) = 1/4$.
Lemmas~\ref{lem:term2_bound} and~\ref{lem:vbar_bound} give
$\Prob(|T_2| > r^2/8) \le 2e^{-K/128} + 2e^{-K/100}$.

\emph{Total.} At most $6e^{-K/2048} \le 3\delta'/4$ exactly as in
Lemma~\ref{lem:test_error_h0}.
\end{proof}

\begin{definition}[Multi-scale algorithm; Algorithm~2]\label{def:multiscale_algorithm}
\lean{Maiti2026Power.nsScale, Maiti2026Power.nsDelta, Maiti2026Power.nsS, Maiti2026Power.nsK, Maiti2026Power.nsJ, Maiti2026Power.nsStart, Maiti2026Power.one_le_nsJ, Maiti2026Power.nsScale_nsJ_lt, Maiti2026Power.NormEstParam.testStat, Maiti2026Power.NormEstParam.testH1, Maiti2026Power.NormEstParam.jStar}\leanok
\uses{def:scale_test}
Let $\eps \in (0,1]$, $\delta \in (0,1)$ and set, for $j \in \N$,
\[
  t_j := 2^j\eps, \qquad \delta_j := \frac{\delta}{2^{j+2}}, \qquad
  s_j := \Big\lceil \frac{4d}{t_j^2} \Big\rceil, \qquad
  K_j := \big\lceil 5120\log(1/\delta_j) \big\rceil, \qquad
  J := \min\{ j \in \N : t_j \ge 2\sqrt d \} .
\]
Since $\eps \le 1 < 2\sqrt d$, $J \ge 1$; and $t_J < 4\sqrt d$ by minimality. The
\emph{multi-scale algorithm} runs the tests $\mathrm{Test}(t_j, \delta_j)$, $j = 0, \dots, J-1$,
one after the other: test $j$ occupies the rounds $[\mathrm{start}_j, \mathrm{start}_{j+1})$
with $\mathrm{start}_j := \sum_{i<j} s_iK_i$, and its budget is the deterministic number
$N_{\mathrm{ms}} := \mathrm{start}_J = \sum_{j<J} s_jK_j$. Its output is the index
\[
  j_* := \min\big( \{ j < J : \mathrm{Test}(t_j,\delta_j) \text{ returns } H_0 \} \cup \{J\} \big),
\]
i.e.\ the first scale at which the test returns $H_0$, or $J$ if all $J$ tests return $H_1$,
and the constant-factor estimate is $r_0 := t_{j_*}$. It satisfies
$\eps \le r_0 \le t_J < 4\sqrt d$ and $r_0 \in \{t_0, \dots, t_J\}$.
\end{definition}

\textbf{Lean remark.} The paper stops the loop at $j_*$ and idles afterwards; since the
budget is deterministic anyway, the formalization runs all $J$ tests (the tests after $j_*$
are simply not used by the output), which makes the sampling rule of the first phase
non-adaptive. The output $j_*$ is a measurable function of the observations of the rounds
$t < N_{\mathrm{ms}}$ (\texttt{NormEstParam.jStar}, defined by \texttt{Nat.find} from the
test outcomes \texttt{testH1}), with finitely many values; the multi-scale algorithm is the
first phase of the meta-algorithm of Definition~\ref{def:meta_algorithm}.

\textbf{Lean remark.} The loop is a sequential composition of $J$ fixed-length phases
(Lemma~\ref{lem:composition}) whose parameters depend on the past: phase $j$ is
$\mathrm{Test}(t_j,\delta_j)$ if no earlier test returned $H_0$ (a measurable function of the
history, computed from the statistics of the earlier blocks) and the constant algorithm
playing $0$ otherwise. The data-dependent output $r_0$ is a measurable function of the full
history of length $N_{\mathrm{ms}}$ with finitely many values. The algorithm is adaptive, but
in a very restricted sense: only the choice between ``continue'' and ``idle'' depends on the
data.

\begin{lemma}[Good event of the multi-scale algorithm]\label{lem:norm_multiscale_good_event}
\lean{Maiti2026Power.NormEstParam.testOK, Maiti2026Power.NormEstParam.badTest, Maiti2026Power.NormEstParam.nsMeasure_real_badTest_le, Maiti2026Power.NormEstParam.bad, Maiti2026Power.NormEstParam.nsMeasure_real_bad_le, Maiti2026Power.sum_nsDelta_le}\leanok
\uses{def:multiscale_algorithm, lem:test_error_h0, lem:test_error_h1, lem:statistic_structure}
Run the multi-scale algorithm against $\nu_\theta$ on the seed space of
Lemma~\ref{lem:statistic_structure}. For $j < J$ let $E_j$ be the event that the outcome of
test $j$ is wrong: it returns $H_1$ while $r \le t_j$, or $H_0$ while $r \ge 2t_j$.
Then $\Prob_{\mathrm{seed}}(E_j) \le 3\delta_j/4$ and
$\Prob_{\mathrm{seed}}\big( \bigcup_{j<J} E_j \big) \le 3\delta/8$.
\end{lemma}

\begin{proof}
\leanok
\uses{lem:statistic_structure, lem:norm_meta_plan, lem:test_error_h0, lem:test_error_h1, def:multiscale_algorithm}
By Lemma~\ref{lem:norm_meta_plan}, the observations of the window of test $j$ are
$y_{\mathrm{start}_j + ks + \ell} = \ip{\sigma_{\mathrm{start}_j+ks}/\sqrt d}{\theta}
+ \eta_{\mathrm{start}_j+ks+\ell}$, so by Lemma~\ref{lem:statistic_structure}(iii) the
statistic $U_j$ of test $j$ is $\bar Z$ evaluated at the window projection
$\pi_{\mathrm{start}_j, s_j, K_j}$, whose law is the block law $\Prob_{s_j,K_j}$
(Lemma~\ref{lem:statistic_structure}(i)). Hence $E_j$ is the preimage under this
projection of the error event of the test on the block space, and its probability is at
most $3\delta_j/4$ by Lemma~\ref{lem:test_error_h0} (if $r \le t_j$) or
Lemma~\ref{lem:test_error_h1} (if $r \ge 2t_j$); when $t_j < r < 2t_j$, $E_j = \emptyset$.
Here $\delta_j \le \delta/4 < 1/4$. The union bound gives
$\sum_{j<J} 3\delta_j/4 \le \frac{3\delta}{4}\sum_{j\ge0}2^{-j-2} = \frac{3\delta}{8}$.
\end{proof}

\begin{theorem}[Guarantee of the multi-scale algorithm]\label{thm:multiscale_guarantee}
\lean{Maiti2026Power.NormEstParam.jStar_bounds}\leanok
\uses{def:multiscale_algorithm, lem:norm_multiscale_good_event}
Let $E_j$, $j < J$, be the events of Lemma~\ref{lem:norm_multiscale_good_event} and let $G$
be the complement of $\bigcup_{j<J}E_j$ (all tests are correct), so
$\Prob_{\mathrm{seed}}(G) \ge 1 - 3\delta/8$ for every $\theta$. On $G$ the output $j_*$
satisfies one of the following:
\begin{enumerate}
  \item[(a)] $j_* = 0$ (that is, $r_0 = \eps$) and $r < 2\eps$;
  \item[(b)] $0 < j_* < J$, $r/2 < t_{j_*} \le 2r$, and $\eps \le t_{j_*} < 2\sqrt d$;
  \item[(c)] $j_* = J$ (that is, $r_0 = t_J \ge 2\sqrt d$) and $r \ge \sqrt d$.
\end{enumerate}
\end{theorem}

\begin{proof}
\leanok
\uses{lem:norm_multiscale_good_event, def:multiscale_algorithm}
The probability bound is Lemma~\ref{lem:norm_multiscale_good_event}. Fix an outcome in $G$:
every test $j < J$ returns $H_1$ if $r \ge 2t_j$ and $H_0$ if $r \le t_j$. Recall
$t_{j+1} = 2t_j$, $J \ge 1$ and $t_J \ge 2\sqrt d$.

If $j_* = 0$: test $0$ returned $H_0$, which is correct only if $r < 2t_0 = 2\eps$.

If $j_* > 0$: test $j_* - 1$ returned $H_1$, which is correct only if $r > t_{j_*-1} = t_{j_*}/2$,
so $t_{j_*} < 2r$. If moreover $j_* < J$, test $j_*$ returned $H_0$, so $r < 2t_{j_*}$, i.e.\
$r/2 < t_{j_*}$; and $\eps \le t_{j_*} < 2\sqrt d$ by the definition of $J$: this is (b). If
$j_* = J$ then $\sqrt d \le t_J/2 = t_{J-1} < r$: this is (c).
\end{proof}

\begin{lemma}[Budget of the multi-scale algorithm]\label{lem:multiscale_samples}
\lean{Maiti2026Power.sum_pow_le, Maiti2026Power.sum_mul_pow_le, Maiti2026Power.NormEstParam.s_le, Maiti2026Power.NormEstParam.K_le, Maiti2026Power.NormEstParam.T₁_le}\leanok
\uses{def:multiscale_algorithm}
With $L := \log(4/\delta)$, the budget of the multi-scale algorithm satisfies
\[
  N_{\mathrm{ms}} = \sum_{j<J} s_jK_j
  \le \frac{8d}{\eps^2}\Big( 6828\,L + 1578 \Big)
  \le 63736\,\frac{d\log(4/\delta)}{\eps^2}
  \le C'\,\frac{d\log(4/\delta)}{\eps^2}, \qquad C' := 64000 .
\]
\end{lemma}

\begin{proof}
\leanok
\uses{def:multiscale_algorithm}
For $j < J$, $t_j < 2\sqrt d$, so $4d/t_j^2 > 1$ and
$s_j \le 4d/t_j^2 + 1 \le 8d/t_j^2 = 8d\,4^{-j}/\eps^2$. Next,
$\log(1/\delta_j) = L + j\log 2$, so
$K_j \le 5120L + 5120\,j\log 2 + 1 \le 5121\,L + 3549\,j$ using $L \ge \log 4 > 1$ and
$5120\log 2 \le 3549$. With $\sum_{j\ge0}4^{-j} = 4/3$ and $\sum_{j\ge0}j4^{-j} = 4/9$,
\[
  N_{\mathrm{ms}} \le \frac{8d}{\eps^2}\sum_{j\ge0}4^{-j}(5121L + 3549j)
  = \frac{8d}{\eps^2}\Big( \frac{4}{3}\,5121\,L + \frac49\,3549 \Big)
  \le \frac{8d}{\eps^2}\big( 6828\,L + 1578 \big).
\]
Finally $1578 \le 1139\,L$ (as $L \ge \log 4 \ge 1.386$), so
$N_{\mathrm{ms}} \le 8\cdot 7967\,dL/\eps^2 \le 64000\,dL/\eps^2$.
\end{proof}

\section{The large-norm regime}
\label{sec:norm_large}

\begin{definition}[Large-norm estimator; Algorithm~3]\label{def:large_norm_algorithm}
\lean{Maiti2026Power.lnN, Maiti2026Power.lnNoise, Maiti2026Power.lnDelta, Maiti2026Power.lnR, Maiti2026Power.lnEst}\leanok
\uses{def:fixed_design, def:norm_estimator}
Let $\eps \in (0,1]$, $\delta \in (0,1)$ and $n := \lceil 48\log(4/\delta)/\eps^2 \rceil$. The
\emph{large-norm estimator} is the fixed-design estimation algorithm with budget $dn$ which
plays each basis vector $e_i \in \ball_d$, $i = 1,\dots,d$, exactly $n$ times (in any fixed
order), computes the empirical means $\hat\theta_i := \frac1n\sum_{t : x_t = e_i} y_t$, and
outputs
\[
  \hat R := \norm{\hat\theta}^2 - \frac dn, \qquad
  \hat r := \sqrt{\max(\hat R, 0)} .
\]
On the \emph{noise space} $(\R^n)^d$ with the law $\Normal(0,1)^{\otimes dn}$ of the noises
$\eta = (\eta_{i,\ell})_{i<d,\ell<n}$ of these rounds, $\hat\theta = \theta + \Delta$ with
$\Delta_i := \frac1n\sum_{\ell<n}\eta_{i,\ell}$, and $\hat R$, $\hat r$ are functions of
$(\theta, \eta)$.
\end{definition}

\textbf{Lean remark.} The basis design is the second phase of the meta-algorithm of
Definition~\ref{def:meta_algorithm} when $j_* = J$ (rounds $T_1 + in + \ell$ play $e_i$);
the guarantee below is stated on the noise space, and transported through the basis window
projection of Lemma~\ref{lem:statistic_structure}(ii).

\begin{lemma}[Decomposition of the large-norm estimator]\label{lem:large_norm_decomposition}
\lean{Maiti2026Power.lnR_sub_eq, Maiti2026Power.hasLaw_two_mul_inner_lnDelta, Maiti2026Power.lnG, Maiti2026Power.hasLaw_lnG, Maiti2026Power.norm_lnDelta_sq_sub_eq, Maiti2026Power.hasSubexponentialMGF_lnQ}\leanok
\uses{def:large_norm_algorithm, def:subexp}
On the noise space, $\Delta \sim \Normal(0, I_d/n)$ and
\[
  \hat R - r^2 = L + Q, \qquad
  L := 2\ip{\theta}{\Delta} \sim \Normal\big(0, 4r^2/n\big), \qquad
  Q := \norm{\Delta}^2 - \frac dn = \frac{\norm{g}^2 - d}{n},
\]
where $g := \sqrt n\,\Delta \sim \Normal(0, I_d)$; $Q$ has
$\mathrm{HasSubexponentialMGF}\ Q\ (8d/n^2)\ (4/n)$.
\end{lemma}

\begin{proof}
\leanok
\uses{lem:gauss_pi, lem:gauss_linear_image, lem:chi_square_subexp, lem:subexp_const_mul}
The coordinate $\hat\theta_i$ is $\theta_i$ plus the average of the $n$ noise variables of
the rounds where $e_i$ is played. These $d$ averages are independent
$\Normal(0,1/n)$ (a sum of independent Gaussians is Gaussian,
Lemmas~\ref{lem:gauss_pi} and~\ref{lem:gauss_linear_image}), so
$\Delta \sim \Normal(0, I_d/n)$. Expanding
$\norm{\theta + \Delta}^2 = r^2 + 2\ip{\theta}{\Delta} + \norm{\Delta}^2$ gives the
decomposition; $L$ is a linear image of $\Delta$ with variance $4\theta^\top(I/n)\theta$
(Lemma~\ref{lem:gauss_linear_image}). Finally $\norm{g}^2 - d$ has parameters $(8d, 4)$ by
Lemma~\ref{lem:chi_square_subexp}, and scaling by $1/n$ gives $(8d/n^2, 4/n)$ by
Lemma~\ref{lem:subexp_const_mul}.
\end{proof}

\begin{theorem}[Guarantee in the large-norm regime]\label{thm:large_norm_guarantee}
\lean{Maiti2026Power.lnNoise_real_lnEst_gt_le, Maiti2026Power.NormEstParam.n_le}\leanok
\uses{def:large_norm_algorithm}
Let $\eps \in (0,1]$, $\delta \in (0,1)$ and $\theta$ with $r \ge \sqrt d$. Then, on the noise
space,
\[
  \Prob\big( \abs{\hat R - r^2} \ge \eps r \big) \le \frac{\delta}{2},
  \qquad\text{hence}\qquad
  \Prob\big( \abs{\hat r - r} \le \eps \big) \ge 1 - \frac{\delta}{2},
\]
and the budget satisfies $n \le 48\log(4/\delta)/\eps^2 + 1 \le 49\log(4/\delta)/\eps^2$,
so $dn \le 49\,d\log(4/\delta)/\eps^2$.
\end{theorem}

\begin{proof}
\leanok
\uses{lem:large_norm_decomposition, lem:gauss_tail, lem:subexp_tail, lem:gauss_law_of_cov}
Write $L_\delta := \log(4/\delta) > 1$; then $n \ge 48L_\delta/\eps^2 \ge 48$ and
$n\eps^2 \ge 48L_\delta$. By Lemma~\ref{lem:large_norm_decomposition},
$\{|\hat R - r^2| \ge \eps r\} \subseteq \{|L| \ge \eps r/2\} \cup \{|Q| \ge \eps r/2\}$.

\emph{Gaussian term.} $L \sim \Normal(0, 4r^2/n)$ and $-L$ has the same law
(Lemma~\ref{lem:gauss_law_of_cov}), so by Lemma~\ref{lem:gauss_tail} applied to $L$ and $-L$,
\[
  \Prob(|L| \ge \eps r/2) \le 2\exp\Big( -\frac{(\eps r/2)^2}{2\cdot 4r^2/n} \Big)
  = 2\exp\Big( -\frac{n\eps^2}{32} \Big)
  \le 2\exp\Big( -\frac{3}{2}L_\delta \Big) = 2\Big(\frac{\delta}{4}\Big)^{3/2}
  = \frac{\delta^{3/2}}{4} \le \frac{\delta}{4}.
\]

\emph{Chi-square term.} Lemma~\ref{lem:subexp_tail} (both tails) for $Q$ with parameters
$(V,b) = (8d/n^2, 4/n)$ and $u := \eps r/2$:
$u^2/(2V) = \eps^2r^2n^2/(64d) \ge n^2\eps^2/64$ (as $r^2 \ge d$), and
$n^2\eps^2/64 = n(n\eps^2)/64 \ge 48\cdot 48 L_\delta/64 = 36L_\delta$; and
$u/(2b) = \eps r n/16 \ge \eps n/16 \ge 48L_\delta/(16\eps) \ge 3L_\delta$ (as
$r \ge \sqrt d \ge 1$ and $\eps \le 1$). Hence
$\Prob(|Q| \ge \eps r/2) \le 2\exp(-3L_\delta) = 2(\delta/4)^3 \le \delta/4$.

\emph{Union bound.} $\Prob(|\hat R - r^2| \ge \eps r) \le \delta/2$.

\emph{From $r^2$ to $r$.} On the event $|\hat R - r^2| < \eps r$, since $\eps \le 1 \le r$,
$\hat R > r(r - \eps) \ge 0$, so $\hat r = \sqrt{\hat R}$ and
$|\hat r - r| = |\hat R - r^2|/(\sqrt{\hat R} + r) < \eps r/r = \eps$.

\emph{Budget.} $n \le 48L_\delta/\eps^2 + 1$ and $1 \le L_\delta/\eps^2$.
\end{proof}

\begin{remark}[On the deterministic design and the constant $48$]\label{rem:norm_large_design}
The paper draws $n$ random Rademacher rows and uses a random-matrix concentration bound
(Vershynin \cite{vershynin2018high}, Section~4.7) to control the least-squares covariance
$(X^\top X)^{-1}$, together with the Gaussian Hanson--Wright inequality. With the
deterministic basis design the covariance of $\hat\theta$ is exactly $I_d/n$, so only the
one-dimensional Gaussian tail and the chi-square tail are needed, and no invertibility
event has to be handled. The value $n = \lceil 32\log(4/\delta)/\eps^2\rceil$ (which is what
one would first write down, and what the blueprint outline suggested) is not sufficient
for the stated failure probability $\delta/2$: it makes the Gaussian term alone equal to
$2\exp(-\log(4/\delta)) = \delta/2$, leaving nothing for the chi-square term. With $48$ the
two terms are $\delta^{3/2}/4$ and $\delta^3/32$, whose sum is at most $\delta/2$ for every
$\delta \in (0,1)$. The paper also assumes $\eps \le r/2$
for the final step; with $\eps \le 1 \le \sqrt d \le r$ the weaker $\eps \le r$ suffices.
\end{remark}

\section{The meta-algorithm}
\label{sec:norm_meta}

\begin{definition}[Meta-algorithm for norm estimation; Algorithm~4]\label{def:meta_algorithm}
\lean{Maiti2026Power.NormEstParam, Maiti2026Power.nsN2, Maiti2026Power.nsBudget, Maiti2026Power.NormEstParam.T, Maiti2026Power.NormEstParam.phaseOf, Maiti2026Power.NormEstParam.rbBlock, Maiti2026Power.NormEstParam.basisRound, Maiti2026Power.NormEstParam.planAction, Maiti2026Power.NormEstParam.nextAction, Maiti2026Power.NormEstParam.alg, Maiti2026Power.NormEstParam.addStat, Maiti2026Power.NormEstParam.lnRStat, Maiti2026Power.NormEstParam.estimate, Maiti2026Power.NormEstParam.output}\leanok
\uses{def:multiscale_algorithm, def:additive_algorithm, def:large_norm_algorithm, lem:seeded, def:norm_estimator}
Let $\eps \in (0,1]$, $\delta \in (0,1)$, let $T_1 := N_{\mathrm{ms}}$ be the budget of the
multi-scale algorithm and
\[
  N_2 := \max\big( dn,\ \max\{ s(t_j)K(t_j) : j < J \} \big),
\]
where $s(t) = \lceil 4d/t^2\rceil$, $K(t) = \lceil 4096\,t^2\log(4/\delta)/\eps^2\rceil$ are the
parameters of Algorithm~1 with $r_0 = t$ (Definition~\ref{def:additive_algorithm}) and $n$
is that of Definition~\ref{def:large_norm_algorithm}. The \emph{meta-algorithm}
$\mathrm{NormEst}(\eps,\delta)$ is the seeded algorithm (Lemma~\ref{lem:seeded}) with budget
$N_{\mathrm{norm}}(\eps,\delta) := T_1 + N_2$ and seeds uniform on $\{-1,+1\}^d$, defined by the
following \emph{round schedule}. The first phase is the multi-scale algorithm of
Definition~\ref{def:multiscale_algorithm}: for $t < T_1$, if $t = \mathrm{start}_j + ks + \ell$
with $j < J$, $k < K_j$, $\ell < s_j$, the round $t$ belongs to the \emph{Rademacher block}
starting at $b := \mathrm{start}_j + ks$, and the action is $x_t := \sigma_b/\sqrt d$ where
$\sigma_b$ is the seed of round $b$ (the action of round $b$, repeated). The second phase,
rounds $T_1 \le t < T_1 + N_2$, depends on the output $j_*$ of the first phase:
\begin{enumerate}
  \item[(a)] if $j_* = 0$: the action is $0$ and the estimate is $\hat r := \eps$;
  \item[(b)] if $j_* = J$: the rounds $T_1 + in + \ell$ ($i < d$, $\ell < n$) are
        \emph{basis rounds} playing $e_i$, and the estimate is the large-norm estimate
        $\hat r = \sqrt{\max(\hat R, 0)}$ of Definition~\ref{def:large_norm_algorithm} computed
        from their observations;
  \item[(c)] if $0 < j_* < J$: the rounds $T_1 + ks + \ell$ ($k < K(t_{j_*})$,
        $\ell < s(t_{j_*})$) form the Rademacher blocks of Algorithm~1 with $r_0 = t_{j_*}$
        (Definition~\ref{def:additive_algorithm}), and the estimate is
        $\hat r = \sqrt{\max(\bar Z, 0)}$ computed from their observations.
\end{enumerate}
The remaining rounds of the second phase play the action $0$ (their observations are not
used). The \emph{plan} $\mathrm{plan}(j_*, t, \sigma)$ is the action of round $t$ as a
function of $j_*$ and of the seeds $\sigma$ at the block starts; as a seeded algorithm, the
action of round $t$ is computed from the past actions and observations (which determine
$j_*$ once $t \ge T_1$, and the action of the current block start) and the fresh seed of
round $t$.
\end{definition}

\textbf{Lean remark.} The parameters are gathered in a structure \texttt{NormEstParam}
(\texttt{ε}, \texttt{δ} and their ranges). The schedule is encoded by two functions of
$(j_*, t)$: \texttt{rbBlock} returns the start of the Rademacher block of round $t$ (if any)
and \texttt{basisRound} the basis vector of round $t$ (if any); \texttt{planAction} is the
plan and \texttt{nextAction} its computation from the past, so that the seeded algorithm
\texttt{alg} has \texttt{next}$_n(h, u) = \mathrm{nextAction}(n, \mathrm{past}(h), j_*(h), u)$,
where $h$ is the history of the first $n$ rounds.
The estimate is the measurable function \texttt{estimate} of the observation sequence
(\texttt{output} on histories of length $N_{\mathrm{norm}}$), and the identification
algorithm is the fixed-budget algorithm with this deterministic output. The paper composes
the phases sequentially; here the composition is replaced by the deterministic schedule,
and the data-dependence is entirely in $j_*$.

\begin{lemma}[The run of the meta-algorithm follows its plan]\label{lem:norm_meta_plan}
\lean{Maiti2026Power.NormEstParam.seedAction_eq_planAction, Maiti2026Power.NormEstParam.yω_eq, Maiti2026Power.NormEstParam.testStat_eq_rbStat, Maiti2026Power.NormEstParam.addStat_eq_rbStat, Maiti2026Power.NormEstParam.lnRStat_eq_lnR, Maiti2026Power.NormEstParam.estimate_congr, Maiti2026Power.NormEstParam.estimate_yOf_seedFinHist, Maiti2026Power.NormEstParam.jsω_prefixExtend}\leanok
\uses{def:meta_algorithm, lem:seeded, lem:statistic_structure}
Run $\mathrm{NormEst}(\eps,\delta)$ against $\nu_\theta$ on the seed space
$(\Omega_{\mathrm{seed}}, \Prob_{\mathrm{seed}})$ of Lemma~\ref{lem:statistic_structure}, with
seeds $\sigma_t$ and noises $\eta_t$, and let $j_* = j_*(\omega)$ be the output of the first
phase. Then for every round $t$ the action is the plan action,
$x_t = \mathrm{plan}(j_*, t, \sigma)$, and $y_t = \ip{x_t}{\theta} + \eta_t$. Consequently:
\begin{enumerate}
  \item[(i)] for $j < J$, the statistic $U_j$ of test $j$ is $\bar Z$ evaluated at the window
    projection $\pi_{\mathrm{start}_j, s_j, K_j}(\omega)$;
  \item[(ii)] on $\{0 < j_* < J\}$, the statistic $\bar Z$ of the second phase is $\bar Z$
    evaluated at $\pi_{T_1, s(t_{j_*}), K(t_{j_*})}(\omega)$;
  \item[(iii)] on $\{j_* = J\}$, the statistic $\hat R$ of the second phase is the large-norm
    statistic evaluated at the basis window projection $\pi'_{T_1, n}(\omega)$;
  \item[(iv)] $j_*$ is a function of the prefix $(\omega_t)_{t<T_1}$, and the estimate is a
    function of the observations of the rounds $t < N_{\mathrm{norm}}$ only (so the estimate
    of the run is the estimate computed from the history $\hist_{N_{\mathrm{norm}}}$).
\end{enumerate}
\end{lemma}

\begin{proof}
\leanok
\uses{def:meta_algorithm, lem:seeded, lem:statistic_structure}
Strong induction on $t$. For $t < T_1$ the schedule does not depend on $j_*$, and the
action computed from the past is the plan action because the action of the block start
$b \le t$ is, by the induction hypothesis, $\sigma_b/\sqrt d$ (or $t = b$ and the fresh seed is
used). For $t \ge T_1$, the value of $j_*$ computed from the past observations
$(y_r)_{r<T_1}$ is the true $j_*$ (the first phase is complete), and the same argument
applies. The observation identity is the definition of the noise environment.
Parts (i)--(iii) follow from Lemma~\ref{lem:statistic_structure}(iii) and the
corresponding identities for the basis rounds ($\ip{e_i}{\theta} = \theta_i$); part (iv) from
the fact that $j_*$ only reads the rounds $t < T_1$, each of which only depends on the
coordinates $\omega_r$, $r \le t$, and that each statistic of the second phase only reads
rounds $t < T_1 + N_2$.
\end{proof}

\textbf{Lean remark.} The second phase is ``$\mathrm{alg}_2(h)$'' in the notation of
Lemma~\ref{lem:composition}, a measurable function of the first-phase history $h$ through
$r_0(h) \in \{t_0,\dots,t_J\}$; since $r_0$ takes finitely many values, measurability is a
finite case distinction. The estimate $\hat r$ is a deterministic measurable function of the
whole history with values in $[0,\infty)$, i.e.\ a recommendation kernel in the sense of
Definition~\ref{def:norm_estimator}. Padding with a constant action keeps the budget
deterministic, as required by Definition~\ref{def:bai_alg}; the padded rounds are simply
not used by the estimate.

\begin{theorem}[Norm estimation; Theorem~8 of the paper]\label{thm:norm_estimation}
\lean{Maiti2026Power.NormEstParam.fail2, Maiti2026Power.NormEstParam.fail_subset, Maiti2026Power.NormEstParam.nsMeasure_real_jsω_inter_preimage_eq_mul, Maiti2026Power.NormEstParam.nsMeasure_real_fail2_le, Maiti2026Power.NormEstParam.nsMeasure_real_estimate_fail_le, Maiti2026Power.NormEstParam.one_sub_le_seedMeasure_real_output, Maiti2026Power.NormEstParam.N₂_le, Maiti2026Power.NormEstParam.T_le}\leanok
\uses{def:meta_algorithm, def:norm_estimator, lem:seeded}
Let $\eps \in (0,1]$ and $\delta \in (0,1)$. The meta-algorithm $\mathrm{NormEst}(\eps,\delta)$
is an $(\eps,\delta)$-accurate estimation algorithm on $\ball_d$: for every
$\theta \in \R^d$, $\Prob_\theta(|\hat r - \norm{\theta}| \le \eps) \ge 1 - 7\delta/8$. Its
budget satisfies
\[
  N_{\mathrm{norm}}(\eps,\delta) \le C_{\mathrm{norm}}\,\frac{d\log(4/\delta)}{\eps^2},
  \qquad C_{\mathrm{norm}} := 10^5 .
\]
\end{theorem}

\begin{corollary}[Theorem~8 of the paper, existential form]\label{cor:norm_estimation_exists}
\lean{Maiti2026Power.exists_isAccurateNormEst}\leanok
\uses{thm:norm_estimation, def:norm_estimator}
Let $\eps \in (0,1]$ and $\delta \in (0,1)$. There exist a budget
$T \le C_{\mathrm{norm}}\, d\log(4/\delta)/\eps^2$ and an $(\eps,\delta)$-accurate estimation
algorithm on $\ball_d$ with budget $T$.
\end{corollary}
\begin{proof}
\leanok
\uses{thm:norm_estimation, lem:seeded}
Take $\mathrm{NormEst}(\eps,\delta)$ of Theorem~\ref{thm:norm_estimation}; it is
$(\eps, 7\delta/8)$-accurate, hence $(\eps,\delta)$-accurate. (For $d = 0$ every reward
vector is $0$ and the constant estimate $\eps$ with budget $0$ is accurate.)
\end{proof}

\begin{proof}
\leanok
\uses{thm:multiscale_guarantee, thm:additive_estimation, thm:large_norm_guarantee, lem:multiscale_samples, lem:norm_meta_plan, lem:statistic_structure, lem:norm_multiscale_good_event, lem:seeded}
\emph{Correctness.} By Lemma~\ref{lem:seeded} it suffices to prove the bound on the seed
space, where, by Lemma~\ref{lem:norm_meta_plan}(iv), the estimate of the run is the
estimate $\hat r(\omega)$ computed from the observation sequence. Let $G$ be the good event
of Theorem~\ref{thm:multiscale_guarantee}, $\Prob_{\mathrm{seed}}(G^c) \le 3\delta/8$. For
$1 \le j \le J$ let $F_j$ be the failure event of the second phase at outcome $j$: for
$j = J$, $F_J := \{ \eps < |\hat r' - r| \}$ where $\hat r'$ is the large-norm estimate
evaluated at the basis window projection $\pi'_{T_1,n}$; for $j < J$,
$F_j := \{ \eps/2 < |\hat r_j - r| \}$ where $\hat r_j$ is the estimate of Algorithm~1 with
$r_0 = t_j$ evaluated at $\pi_{T_1, s(t_j), K(t_j)}$. By Lemma~\ref{lem:norm_meta_plan}(ii)--(iii),
\[
  \{ \eps < |\hat r - r| \} \subseteq G^c \cup \bigcup_{1 \le j \le J}
    \big( G^c{}^c \cap \{j_* = j\} \cap F_j \big):
\]
on $G \cap \{j_* = 0\}$, $r < 2\eps$ by Theorem~\ref{thm:multiscale_guarantee}(a) and
$\hat r = \eps$, so $|\hat r - r| \le \eps$; on $\{j_* = j\}$ with $j \ge 1$ the estimate is
$\hat r_j$ or $\hat r'$. For each $j \ge 1$: if $G \cap \{j_* = j\} = \emptyset$ the term
vanishes; otherwise Theorem~\ref{thm:multiscale_guarantee}(b)--(c) gives the hypotheses of
Theorem~\ref{thm:additive_estimation} ($\eps \le t_j$, $r/2 < t_j \le 2r$) or of
Theorem~\ref{thm:large_norm_guarantee} ($r \ge \sqrt d$), and
$\Prob_{\mathrm{seed}}(F_j) \le 3\delta/8 \le \delta/2$ (resp.\ $\le \delta/2$) by
Lemma~\ref{lem:statistic_structure}(i)--(ii) (the law of the window projection is the
block, resp.\ noise, law). Since $\{j_* = j\}$ is a function of the prefix
$(\omega_t)_{t<T_1}$ (Lemma~\ref{lem:norm_meta_plan}(iv)) and $F_j$ is a function of a window
projection with $a = T_1$, they are independent (Lemma~\ref{lem:statistic_structure}), so
$\Prob_{\mathrm{seed}}(\{j_* = j\} \cap F_j) = \Prob_{\mathrm{seed}}(j_* = j)\,\Prob_{\mathrm{seed}}(F_j)
\le \Prob_{\mathrm{seed}}(j_* = j)\,\delta/2$. Summing over $j$ (the events $\{j_* = j\}$ are
disjoint, so their probabilities add up to at most $1$),
$\Prob_{\mathrm{seed}}(\eps < |\hat r - r|) \le 3\delta/8 + \delta/2 = 7\delta/8$.

\emph{Budget.} Lemma~\ref{lem:multiscale_samples} gives $N_{\mathrm{ms}} \le 63736\,dL/\eps^2$
with $L = \log(4/\delta)$. For $j < J$, $t_j \in [\eps, 2\sqrt d)$, so
Theorem~\ref{thm:additive_estimation} (budget clause with $r_0 = t_j$) gives
$s(t_j)K(t_j) \le 32776\,dL/\eps^2$, and Theorem~\ref{thm:large_norm_guarantee} gives
$dn \le 49\,dL/\eps^2$; hence $N_2 \le 32776\,dL/\eps^2$ and
$N_{\mathrm{norm}} \le (63736 + 32776)\,dL/\eps^2 \le 10^5\,dL/\eps^2$.
\end{proof}

\begin{remark}[Reading of the constants]\label{rem:norm_constants}
The budget is $O(d\log(1/\delta)/\eps^2)$ as in the paper: for $\delta \le 1/4$,
$\log(4/\delta) \le 2\log(1/\delta)$. The constant $C_{\mathrm{norm}} = 10^5$ is dominated by the
multi-scale phase ($63736$), itself driven by the test constant $c_1' = 5120$; the
constants are what the union bounds produce and were not optimized. In the paper the
sample complexity of the multi-scale phase is written $\sum_j s_jt_j$, a typo for
$\sum_j s_jK_j$. The paper's failure probabilities are $\delta/2$ for the multi-scale phase
and $\delta/2$ for the second phase; the sharper $3\delta/8$ of
Lemma~\ref{lem:norm_multiscale_good_event} is the paper's own computation
$\sum_j 3\delta_j/4$.
\end{remark}
